\section{Conclusion}

We presented \ours{}, a training-free concept-erasure framework for text-to-image diffusion models structured as three decoupled stages---\emph{when} (CAS gate), \emph{where} (family-indexed dual-probe winner-take-all), and \emph{how} (anchor inpainting or hybrid guidance)---all driven by a small, user-editable exemplar pack. Across four adversarial nudity benchmarks and four MJA categories, \ours{} substantially outperforms all training-free baselines---including the strongest published compounds (SAFREE~$+$~Safe\_Denoiser, SAFREE~$+$~SGF)---while preserving COCO image quality and transferring to SD~3 (MMDiT) and FLUX (DiT) without design changes. The family decomposition and the pixel-wise winner-take-all partition are the two mechanisms that let a single concept direction split into semantically independent sub-directions and route each detected pixel to the appropriate one, which is what makes simultaneous multi-concept erasure tractable inside the same update.

\paragraph{Limitations.}
(i)~Highly sophisticated adversarial prompts (MMA-Diffusion) remain an active challenge: performance has improved substantially over prior work but a non-trivial fraction still evades all training-free defenses we evaluated. (ii)~Simultaneous multi-concept erasure shows measurable degradation over single-concept erasure, particularly for semantically overlapping concepts (\eg, nudity~$+$~hate); a more principled treatment of cross-concept interference is left for future work. (iii)~The CAS threshold $\tau$ requires per-backbone retuning; we find $\tau{\in}[0.4,0.6]$ works well for $\epsilon$-prediction UNets and $\tau{\in}[0.2,0.4]$ is needed for the velocity-parameterization SD~3. (iv)~On a small number of image-probe-dominant concepts whose paired anchors are semantically ambiguous (MJA-Illegal), SAFREE's global token projection remains competitive with our spatially selective approach.

\paragraph{Broader impact.}
\ours{} is designed to mitigate safety concerns in generative models and contributes to responsible AI deployment. The same spatial localization techniques could in principle be repurposed to amplify harmful content if applied maliciously, and should be deployed as one component of a comprehensive safety pipeline that includes upstream prompt filtering and downstream content review.
